\documentclass{beamer}
\mode<presentation> {\usetheme{default}}
%\setbeamersize{text margin left=10mm,text margin right=10mm} 

\setbeamertemplate{footline}[frame number]

\usepackage{multirow}
\usepackage{amssymb}
\usepackage{amsthm}
\usepackage{amsmath}
\usepackage{booktabs}
\usepackage{bm}
\usepackage{needspace}
\usepackage{comment}
\usepackage{tikz}
\usepackage{graphicx}
\usepackage{bbm, dsfont}
\usepackage{listings}

\usetikzlibrary{arrows,shapes,calc,positioning}

\theoremstyle{plain}
\newtheorem{thm}{Theorem}[section]
\newtheorem{cor}[thm]{Corollary}
\newtheorem{lem}[thm]{Lemma}
\newtheorem{prop}[thm]{Proposition}
%
\theoremstyle{definition}
\newtheorem{defn}{Definition}[section]
\newtheorem{remk}{Definition}[section]
\newtheorem{exam}{Example}[section]
\newtheorem{prob}{Exercise}[section]
\newtheorem{solu2}{Solution}[section]
\newenvironment{Shaded}{\begin{snugshade}}{\end{snugshade}}

\newcommand{\dif}{\operatorname{d}}
\newcommand{\var}{\operatorname{Var}}
\newcommand{\brtexttt}[1]{\texttt{\color{brown}{#1}}}

%opening
\title{Logistic Regression}
\author{Himchan Jeong, FSA}
\date{}

\begin{document}
\maketitle

\section{Logistic Regression}

\subsection{Introduction}


\begin{frame}{Introduction and acknowledgement}

\begin{itemize}
\item I want to start with a disclaimer in the name `logistic \textbf{regression}', which can be in fact misleading.
\item Logistic regression is basically one of the methods for \textbf{binary classification problems}, where the response variable is in fact qualitative and has only two types of responses.
\item Throughout this section, we assume that the response variable $y$ can take only two values, $0$ or $1$, while their interpretation solely depends on the context.
\item Note that most of the contents in this section are from Section 4.3 of the  \href{https://hastie.su.domains/ISLR2/ISLRv2_corrected_June_2023.pdf.download.html}{ISL textbook}.
\end{itemize}

\end{frame}


\begin{frame}{Default data description}

To introduce logistic regression, let us consider a simulated dataset named `Default' with 10,000 observations, where the response variable is \brtexttt{default} (binary) and three covariates are available; \brtexttt{student} (binary), \brtexttt{balance} (continuous), and \brtexttt{income} (continuous).
\begin{table}
\centering
\begin{tabular}{ccccc}
\toprule 
index & \brtexttt{default} & \brtexttt{student} & \brtexttt{balance} & \brtexttt{income} \\
\midrule 
1               & No      & No     & 729.5265  & 44361.625 \\
2               & No      & Yes    & 817.1804  & 12106.135 \\
3               & No      & No     & 1073.5492 & 31767.139 \\
$\vdots$               & $\vdots$      & $\vdots$    & $\vdots$  & $\vdots$ \\
10000  &    No  &   Yes & 200.9222 & 16862.95 \\
\bottomrule
\end{tabular}
\end{table}
Note that 9,000 observations are used to train the models, while 1,000 observations are set aside for the out-of-sample validation.

\end{frame}

\begin{frame}[allowframebreaks]
\frametitle{Simple logistic regression}


As of now, let us only consider \brtexttt{balance} for the logistic regression model.
\begin{itemize}
	\item For simplicity, we write $Pr(\brtexttt{default = Yes} | \brtexttt{balance})$ as $p( \brtexttt{balance})$, which will range between 0 and 1.
	\item Then for any given value of \brtexttt{balance}, a prediction of \brtexttt{default} probability can be made.
	\item Note that the predicted probability needs to be converted to actual decision (default or not) with a threshold. Note that a usual threshold for the binary classification is 0.5 (i.e., if $p( \brtexttt{balance})>0.5$, then that person is expected to default).
	\item However, if the company wants to apply conservative underwriting criteria, then the threshold can be 0.1 (i.e., if $p( \brtexttt{balance})>0.1$, then that person is expected to default and likely denied for a credit application).
\end{itemize}

\framebreak


\begin{itemize}
\item In the language of the GLMs, $y \in \{0, 1\}$ is  \brtexttt{default} (1 is yes), $\mathbf{x}$ is $(1, \brtexttt{balance})'$, and $\beta=(\beta_0, \beta_1)$ is the associated regression coefficients.
\item By using the logit link to connect $p(\mathbf{x})$ and $\mathbf{x}'\beta$, we have
$$
p(\mathbf{x}) = \frac{e^{\mathbf{x}'\beta}}{1+ e^{\mathbf{x}'\beta}} \ \text{ or } \ \log\left(\frac{p(\mathbf{x})}{1-p(\mathbf{x})}\right) = \mathbf{x}'\beta.
$$
\item Note that one can use a different link function to connect $p(\mathbf{x})$ and $\mathbf{x}'\beta$, for example, $\eta(\cdot) = \Phi^{-1}(\cdot)$ (probit link) and it is called `probit regression', where $\Phi(\cdot)$ is the CDF of a standard normal random variable.
\end{itemize}

\framebreak

\begin{itemize}
\item From the regression results (via \texttt{glm} in \texttt{R}), it is shown that $\hat{\beta}_1 =0.0055$; this indicates that an increase in balance is associated with an
increase in the probability of default. 
\item To be precise, a one-unit increase in balance is associated with an increase in the log odds of default by 0.0055 units.
\end{itemize}

\begin{table}
\caption{Simple logistic regression for \brtexttt{default} only with \brtexttt{balance}}
\begin{tabular}{lrrrr}
\toprule
  & Coefficient & Std. error & $z$-statistic & $p$-value\\
\midrule
\brtexttt{Intercept} & -10.5059 & 0.3758 & -28.0 & $<0.0001$\\
\brtexttt{balance} & 0.0054 & 0.0002 & 23.5 & $<0.0001$\\
\bottomrule
\end{tabular}
\end{table}

\vspace{8pt}

Question: how large is the predicted default probability for an applicant whose \brtexttt{balance} is 2,000 with the above results?

\framebreak


\begin{itemize}
\item One can repeat a similar analysis (simple logistic regression) where $\mathbf{x}$ is $(1, \brtexttt{student[Yes]})'$.
\item It indicates that students tend to have higher default probabilities than non-students:
\end{itemize}

\begin{table}
\caption{Simple logistic regression for \brtexttt{default} only with \brtexttt{student}}  \label{tab:2} \vspace{-10pt}
\begin{tabular}{lrrrr}
\toprule
  & Coefficient & Std. error & $z$-statistic & $p$-value\\
\midrule
\brtexttt{Intercept} & -3.5761 & 0.0771 & -46.4 & $<0.0001$\\
\brtexttt{student[Yes]} & {\color{blue}\textbf{0.5296}} & 0.1211 & 4.4 & $<0.0001$\\
\bottomrule
\end{tabular}
\end{table}

\end{frame}

\begin{frame}[allowframebreaks]
\frametitle{Multiple logistic regression}

Indeed, it would be more natural to consider all the available covariates simultaneously so that $\mathbf{x}=(1, x_1, x_2, x_3)$ is $(1, \brtexttt{balance}, \brtexttt{income}$ (in thousands of dollars), $\brtexttt{student[Yes]})'$, $\beta = (\beta_0, \beta_1, \beta_2, \beta_3)'$, and
$$
\log\left(\frac{p(\mathbf{x})}{1-p(\mathbf{x})}\right) = \beta_0+\beta_1 x_1+\beta_2 x_2+\beta_3 x_3.
$$

\begin{table}
\caption{Multiple logistic regression for \brtexttt{default} with all covariates} \label{tab:3} \vspace{-10pt}
\begin{tabular}{lrrrr}
\toprule
  & Coefficient & Std. error & $z$-statistic & $p$-value\\
\midrule
\brtexttt{Intercept} & -10.7161 & 0.5186 & -20.7 & $<0.0001$\\
\brtexttt{balance} & 0.0056 & 0.0002 & 23.2 & $<0.0001$\\
\brtexttt{income} & 0.0034 & 0.0088 & 0.4 & 0.7019\\
\brtexttt{student[Yes]} & {\color{red}\textbf{-0.5106}} & 0.2505 & -2.0 & 0.0415\\
\bottomrule
\end{tabular}
\end{table}

\framebreak
\begin{itemize}
\item The $p$-value associated with the student status is very small, indicating that it is associated with the probability of default. 
\item However, the coefficient for the student variable is {\color{red}\textbf{negative}}, indicating that students are less likely to default than non-students.
\item In contrast, the coefficient for the student variable is {\color{blue}\textbf{positive}} in
Table  \ref{tab:2}. How is it possible for student status to be associated with an
increase in probability of default in Table  \ref{tab:2} and a decrease in probability
of default in Table \ref{tab:3}? 
\end{itemize}

\framebreak

\begin{itemize}
\item For both students and non-students group, higher credit balance is associated with higher default rate.
\item Students tend to have higher credit balance than non-students.
\end{itemize}

\begin{figure}
\vspace{-10pt}
  \centering
\includegraphics[width=0.9\textwidth]{default.png}
\end{figure}
In conclusion, it is shown from the multiple logistic regression model that a student is less risky than a non-student \emph{with the same credit card balance}.


\end{frame}


\begin{frame}[allowframebreaks]
\frametitle{Model evaluation}
\begin{itemize}
	\item The predicted binary class does not only depend on the fitted model, but also the threshold for the decision (e.g., 0.5 vs 0.1) as mentioned before.
	\item As such, one can expect different out-of-sample validation results depending on the used threshold even with the same model used.
	\item Once we have the actual and predicted binary results, there are in total four combinations and they can be represented in the following format - which is called a \textbf{confusion matrix}.

\begin{table}[]
\begin{tabular}{cc|ccc}
                                                                    &       & \multicolumn{3}{c}{Actual}     \\
                                                                             &       & Negative  & \multicolumn{1}{c|}{Positive} & Total \\ \hline
\multirow{2}{*}{\begin{tabular}[c]{@{}c@{}}Predicted\\ \end{tabular}} & Negative    & TN & \multicolumn{1}{c|}{FN}  & TN+FN   \\
                                                                             & Positive   & FP   & \multicolumn{1}{c|}{TP}  & FP+TP    \\ \cline{2-5} 
                                                                             & Total & TN+FP  & \multicolumn{1}{c|}{FN+TP}  &  
\end{tabular}
\end{table}

\framebreak

	\item Here we define two concepts. \textbf{Sensitivity} refers to the model's ability to correctly detect those who are \emph{positive} out of those \emph{who do have} the condition, which is  defined as follows:
$$
\text{Sensitivity:=}\frac{\text{TP}}{\text{TP}+\text{FN}}.
$$
	\item \textbf{Specificity} refers to the model's ability to correctly detect those who are \emph{negative} out of those \emph{who do not have} the condition, which is  defined as follows:
$$
\text{Specificity:=}\frac{\text{TN}}{\text{TN}+\text{FP}}.
$$
\end{itemize}

\framebreak


\begin{table}[]
\caption{Confusion matrix where the threshold is 0.5}
\begin{tabular}{cc|ccc}
                                                                    &       & \multicolumn{3}{c}{Actual default}     \\
                                                                             &       & No  & \multicolumn{1}{c|}{Yes} & Total \\ \hline
\multirow{2}{*}{\begin{tabular}[c]{@{}c@{}}Predicted\\ Default\end{tabular}} & No    & 956 & \multicolumn{1}{c|}{25}  & 981   \\
                                                                             & Yes   & 4   & \multicolumn{1}{c|}{15}  & 19    \\ \cline{2-5} 
                                                                             & Total & 960 & \multicolumn{1}{c|}{40}  & 1000 
\end{tabular}
\end{table}
Sensitivity$=15/40=37.5\%$,   Specificity$=956/960=99.58\%$. 
\vspace{-6pt}
\begin{table}[]
\caption{Confusion matrix where the threshold is 0.1}
\begin{tabular}{cc|ccc}
                           &       & \multicolumn{3}{c}{Actual default}     \\
                                                                             &       & No  & \multicolumn{1}{c|}{Yes} & Total \\ \hline
\multirow{2}{*}{\begin{tabular}[c]{@{}c@{}}Predicted\\ Default\end{tabular}} & No    & 908 & \multicolumn{1}{c|}{5}  & 913   \\
                                                                             & Yes   & 52   & \multicolumn{1}{c|}{35}  & 87    \\ \cline{2-5} 
                                                                             & Total & 960 & \multicolumn{1}{c|}{40}  & 1000 
\end{tabular}
\end{table}
Sensitivity$=35/40=87.5\%$,   Specificity$=908/960=94.58\%$. 

\framebreak

\begin{itemize}
	\item Lower threshold implies that we are more likely to diagnose one as `positive' given the  model. Therefore, the ability to detect `true positive' (which is \textbf{sensitivitiy}) increases in this case while the ability to detect `true negative' (which is \textbf{specificity}) decreases.
	\item On the other hand, a higher threshold implies that we are more likely to diagnose one as `negative' given the model. Therefore, the ability to detect `true negative' (which is \textbf{specificity}) increases in this case while the ability to detect `true positive' (which is \textbf{sensitivitiy}) decreases.
\end{itemize}

\framebreak

The ROC (receiver operating characteristics) curve is a popular graphic for simultaneously displaying \textbf{sensitivitiy} and \textbf{specificity} for all possible thresholds. 

\begin{figure}
  \centering
\includegraphics[trim={0 0 0 1cm},clip, width=0.7\textwidth]{rocs.png}
\end{figure}
\begin{itemize}
	\item The overall performance of a classifier, summarized over all possible thresholds, is given by the \emph{area under the curve} (AUC). An ideal ROC curve will hug the top left corner, so the larger AUC the better the classifier.
	\item Note that if one uses a `random classifier' (which assigns 0 or 1 at random), then the ROC will be the diagonal line and the corresponding AUC will be 0.5.
	\item In this case, the AUC of the (multiple) logistic regression model is 0.972, which is quite satisfactory.
\end{itemize}

\end{frame}




\end{document}